% Procedencia íntegra: CONTRAANGULO_ELECTRON_BARBERO_CATALAN.md, §§2--10.
% La gestión, el historial y los localizadores se conservan en gestion/.
\section{Composición y restitución de las lecturas angular, electrónica y constitutiva}
\label{md:compos:seccion}

\subsection{Coeficiente de acción y construcción del contraángulo}
\label{md:compos:contraangulo}

La abreviatura funcional $B(\alpha)$ reúne una evaluación en $\alpha$; su
argumento pertenece a la función y no representa el producto de una constante
independiente por $\alpha$. Para mantener explícitas las operaciones que
intervienen en esta abreviatura, se escribe
\begin{equation}
 B(\alpha)=\frac{9}{16}\alpha^2-\frac{5}{9}\alpha^3
 +\frac{7}{48}\alpha^4-\frac{1}{54}\alpha^5
 -\frac{90}{\pi}\mathcal C_\pi(\alpha),
 \label{md:compos:b-alpha}
\end{equation}
\begin{equation}
 \mathcal C_\pi(\alpha)=169\alpha^6+
 \frac{D_A\alpha^7}{1-D_A\alpha},
 \qquad D_A=\exp\!\left(-\frac{100\pi\alpha}{9}\right).
 \label{md:compos:correccion-regional}
\end{equation}
La identidad completa es
\begin{equation}
 C^*_{\deg}=\sqrt{\frac{1000\alpha}{\varphi}}+2B(\alpha).
 \label{md:compos:contraangulo-abreviado}
\end{equation}
En particular, $\alpha$ pertenece al radicando, y el término radical se
conserva junto a $2B(\alpha)$. Esta identidad reescribe el mapa construido
anteriormente y mantiene su objeto y su genealogía.

La definición original conserva separados el coeficiente adimensional y la
sección de acción:
\begin{equation}
 \eta_{\mathrm{ret}}=H_5-\frac{90}{\pi}\mathcal C_\pi,
 \qquad
 \hbar_{\mathrm{ret}}=\eta_{\mathrm{ret}}\,10^{-34}\mathcal U_S.
 \label{md:compos:seccion-accion}
\end{equation}
Por tanto,
\begin{equation}
 \boxed{C^*_{\deg}=2\left(
 \frac{\hbar_{\mathrm{ret}}}{10^{-34}\mathcal U_S}+\alpha\right)
 =2(\eta_{\mathrm{ret}}+\alpha).}
 \label{md:compos:contraangulo-accion}
\end{equation}
La expresión utiliza el coeficiente adimensional de la sección de Planck
reducida. La constante $h$ se obtiene después mediante $h=2\pi\hbar$.
La composición consiste en extraer el coeficiente de acción, componerlo con
$\alpha$ y aplicar la clausura bidireccional. Así se mantiene el tipo
dimensional de cada operación.

La construcción de masas expresa partícula y antipartícula como dos
secciones orientadas de una misma banda genealógica. La operación
dodecafásica es
\begin{equation}
 C_M(\tau)=-\operatorname{rev}(\tau),
 \qquad \tau_e=+++---+++---,
 \qquad C_M(\tau_e)=\tau_e,
 \label{md:compos:conjugacion-dodecafasica}
\end{equation}
prolongada al registro y a la representación de carga. La ruta central es
fija bajo $C_M$, mientras la representación de carga distingue electrón y
positrón. Esta composición da contenido preciso a la lectura bidireccional.
El factor dos angular conserva su función de clausura; la conjugación física
conserva además los operadores indicados. El retorno a $216$ restituye el
levantamiento orientado: su estructura no consiste en asignar $108$
actualizaciones a la partícula y otras $108$ a la antipartícula.

\subsubsection{Restitución radial del coeficiente de acción}
\label{md:compos:quinientos}

Escribimos $q_{\mathrm{rad}}=q_\varepsilon$ y $\epsilon=\varepsilon$ para
la razón radial y la orientación de \eqref{rec:eq:eta-orientada}.
El invariante radial orientado $\chi_{\mathrm{rad}}=\epsilon\log q_{\mathrm{rad}}$
permite escribir la restitución previamente construida como
\begin{equation}
 \frac{\eta_{\mathrm{ret}}}{\alpha}
 =-500\tanh\!\left(\frac{\chi_{\mathrm{rad}}}{2}\right)-1.
 \label{md:compos:restitucion-radial}
\end{equation}
Con $\xi=C^*_{\deg}/A_{\deg}$ y $A_{\deg}=1000\alpha$, se obtiene también
\begin{equation}
 \eta_{\mathrm{ret}}=\alpha(500\xi-1).
 \label{md:compos:restitucion-xi}
\end{equation}
Son recuperadores posteriores de la misma sección; su construcción permanece
en $H_5$ y en el retorno regional.

\subsection{Coordenatizaciones angulares y conservación de la acción}
\label{md:compos:cartas}

La coordenatización en grados tiene una vuelta de $360$ unidades y conserva las
particiones $12\cdot30=4\cdot90=3\cdot120=360$. El par angular construido es
\begin{equation}
 A_{\deg}=1000\alpha,
 \qquad C^*_{\deg}=2(\eta_{\mathrm{ret}}+\alpha).
 \label{md:compos:par-angular}
\end{equation}
La conversión se efectúa una vez:
\begin{equation}
 x=\frac{\pi}{180}A_{\deg},\qquad
 y=\frac{\pi}{180}C^*_{\deg},\qquad
 q_\pm=\exp(-x\mp y).
 \label{md:compos:canales}
\end{equation}
Los exponentes utilizan representantes reales levantados y conservan el
número de vueltas. La acción satisface
\begin{equation}
 \hbar\theta_{\mathrm{rad}}
 =\frac{h\theta_{\deg}}{360}.
 \label{md:compos:accion-cartas}
\end{equation}
En el funcional bilineal de Barbero, el transporte de coordenatización da
\begin{equation}
 \frac{\pi A_{\deg}C^*_{\deg}}{180}
 =\frac{180xy}{\pi}.
 \label{md:compos:bilineal-cartas}
\end{equation}
La transformación bilineal forma parte de la misma conversión que se aplica
a los argumentos exponenciales; omitirla cambiaría el funcional.

\subsection{Lectura electrónica completa y retorno de acción}
\label{md:compos:electron}

Se conservan las dos etapas de la misma coordenada electrónica:
\begin{equation}
 \varepsilon_e^{(0)}=\frac{\sqrt{3}}{4}
 \left[\exp\!\left(\frac{\varphi}{\pi^2}\right)-22\alpha^3\right]
 (1+15\Delta_4),\qquad
 \Delta_4=\frac{\pi-e\log\pi}{270}
 \left(1+\frac{\pi}{729}\right),
 \label{md:compos:electron-basal}
\end{equation}
\begin{equation}
 \boxed{\varepsilon_e^\beta=\varepsilon_e^{(0)}
 \left(\frac{H_5}{\eta_{\mathrm{ret}}}\right)^{80-54\alpha+6\Delta_4}.}
 \label{md:compos:electron-completo}
\end{equation}
El centro $(5,5)$, su inversión, las deformaciones $(8,2)/(2,8)$ y el
defecto unitario de cociente son antecedentes del valor. La norma
$\sqrt{3}/4$ procede del conmutador entre proyecciones aditivas y
multiplicativas; el carácter exponencial utiliza la orientación
$\varphi/\pi^2$; $22=27-5$ conserva un complemento regional y $15=3\cdot5$
las incidencias correspondientes. Los lectores de la ruta producen
$(80,54,6)$. El cociente $H_5/\eta_{\mathrm{ret}}$ incorpora el retorno de
acción a la misma sección central.

La función del contraángulo se explicita sustituyendo su inversa demostrada:
\begin{equation}
 \boxed{\varepsilon_e^\beta=\varepsilon_e^{(0)}
 \left(\frac{H_5}{C^*_{\deg}/2-\alpha}\right)^{80-54\alpha+6\Delta_4}.}
 \label{md:compos:electron-contraangulo}
\end{equation}
Esta expresión conserva la generación electrónica anterior. La misma
sección de acción que participa en el contraángulo determina el factor que
lleva la lectura basal a la completa. El cociente cancela la base dimensional
común; la realización absoluta de masa conserva su lector temporal y su coordenatización
dimensional. La sustitución es una composición de dependencias construidas,
sin introducir un ángulo externo como generador electrónico.

\subsection{Funcional de Barbero y transición hexada--octada}
\label{md:compos:barbero}

La incidencia mantiene el mismo par marcado al ampliar el soporte:
\begin{equation}
 \mathcal F_6=\mathsf H\times\binom{\mathsf H}{2}
 \hookrightarrow
 \mathcal F_8=\mathsf O\times\binom{\mathsf H}{2},
 \qquad |\mathcal F_6|=90,\qquad |\mathcal F_8|=120.
 \label{md:compos:incidencia}
\end{equation}
Los cardinales son $6\cdot15$ y $8\cdot15$; el referente compartido permite
comparar ambas incidencias. La diferencia normalizada es $-1/12$. La cadena
dodecafásica conserva doce contribuciones hexádicas y una contribución
octádica de frontera con signo contrario. El lector
\begin{equation}
 S_m(q)=\frac{q^m}{1-q^{3m}}
 =\sum_{j\geq0}q^{m+3mj}
 \label{md:compos:serie-retorno}
\end{equation}
conserva altura y retornos. Así se forma $12S_{90}-S_{120}$ y su diferencia
entre las dos orientaciones. El funcional completo es
\begin{equation}
 \gamma=\frac{\pi A_{\deg}C^*_{\deg}}{180}
 +12\bigl[S_{90}(q_-)-S_{90}(q_+)\bigr]
 -\bigl[S_{120}(q_-)-S_{120}(q_+)\bigr].
 \label{md:compos:funcional-barbero}
\end{equation}
Barbero es, en esta realización, la evaluación orientada conjunta del par
angular y de los retornos de la transición de incidencias. Su participación
en la escala de área se compone después con $\ell_P^2=G\hbar/c^3$ y con el
operador de ocupación. La entropía de entrelazamiento conserva además el
estado reducido y sus probabilidades; su reconstrucción requiere esos datos
y no queda determinada por el área escalar aislada.

\subsubsection{Eliminación de la coordenada común electrón--Barbero}
\label{md:compos:eliminacion}

Definamos $\kappa_e=80-54\alpha+6\Delta_4$. En la rama positiva, con
$\kappa_e\neq0$, la ecuación electrónica implica exactamente
\begin{equation}
 \eta_{\mathrm{ret}}=H_5
 \left(\frac{\varepsilon_e^{(0)}}{\varepsilon_e^\beta}\right)^{1/\kappa_e}.
 \label{md:compos:eliminacion-eta}
\end{equation}
Para expresar la composición manteniendo visibles sus operaciones, pongamos
$x=50\pi\alpha/9$, $y=\pi(\eta_{\mathrm{ret}}+\alpha)/90$ y
$\mathcal R(q)=12S_{90}(q)-S_{120}(q)$. Entonces
\begin{equation}
 \gamma=\frac{100\pi}{9}\alpha(\eta_{\mathrm{ret}}+\alpha)
 +\mathcal R(e^{-x+y})-\mathcal R(e^{-x-y}).
 \label{md:compos:barbero-eta}
\end{equation}
La sustitución de \eqref{md:compos:eliminacion-eta} en esta identidad produce
una relación explícita entre la lectura electrónica completa y Barbero,
conservando $\alpha$, $H_5$, $\Delta_4$ y las incidencias. Es una consecuencia
de compatibilidad de la construcción. La rama física angular utilizada
mantiene $0<y<x$; la operación no introduce una masa metrológica como dato
generador.

En la extensión con etiqueta de orientación $\epsilon=\pm1$,
$C^*_{\deg,\epsilon}=2\epsilon(\eta_{\mathrm{ret}}+\alpha)$, el intercambio
$\epsilon\mapsto-\epsilon$ deja fijo $\eta_{\mathrm{ret}}$, conserva la
expresión electrónica y permuta $q_+$ y $q_-$. El funcional de Barbero cambia
de signo. Esta covariancia distingue una lectura escalar de acción de una
evaluación orientada. La identificación con conjugación de carga conserva
además $C_M$ y la representación electromagnética de hojas; el signo angular
aislado no sustituye esos operadores.

\subsection{Momento orientado de Catalán y respuesta constitutiva}
\label{md:compos:catalan}

El ciclo dodecafásico realiza un plano orientado de cuarto de giro. Su
coeficiente resolvente es $k_4(s)=s/(1+s^2)$. El momento de profundidad dos
queda fijado por
\begin{equation}
 \mathcal C_2(s)=\sum_{k\geq0}
 \frac{(-1)^ks^{2k+1}}{(2k+1)^2},\qquad
 \mathcal D=s\frac{d}{ds},\qquad
 \mathcal D^2\mathcal C_2(s)=k_4(s),\qquad
 \mathcal C_2(1)=G_{\mathrm{Cat}}.
 \label{md:compos:momento-catalan}
\end{equation}
La identidad diferencial se utiliza para $0<s<1$; el valor extremo se
entiende con las condiciones de convergencia y restitución del momento ya
establecidas. La definición intrínseca de Catalán es el valor extremo de ese
momento orientado de profundidad. Su colección conserva más información que
el valor extremo. La respuesta constitutiva utiliza dicha colección:
\begin{equation}
 f(s)=\frac{1+s+s^2}{s(1+s)}\mathcal D^2\mathcal C_2(s)
 =\frac{1-s^3}{1-s^4},\qquad
 r_\pm=f(s_\pm),\qquad s_\pm=q_\pm^{30}.
 \label{md:compos:catalan-respuesta}
\end{equation}
Las potencias $3$ y $4$ corresponden a $90=3\cdot30$ y $120=4\cdot30$.
\begin{proposition}[Restitución funcional con condición de frontera]
\label{md:compos:inversa-catalan}
La colección constitutiva $f$ determina el momento orientado mediante
\[
 \mathcal C_2(s)=\int_0^s\log\frac{s}{t}\,
 \frac{1+t}{1+t+t^2}f(t)\,dt,\qquad 0<s<1.
\]
Es la única solución de
$\mathcal D^2\mathcal C_2(s)=s(1+s)f(s)/(1+s+s^2)$
que tiende a cero en el origen.
\end{proposition}
\begin{proof}
Para la colección construida,
$(1+t)f(t)/(1+t+t^2)=1/(1+t^2)$. La integral converge, y su
derivada logarítmica es
$\mathcal D\mathcal C_2(s)=\int_0^s(1+t^2)^{-1}\,dt$.
Aplicar de nuevo $\mathcal D$ da $s/(1+s^2)$. La integral es
$O(s)$ en el origen. Dos soluciones difieren en $a+b\log s$;
el límite nulo obliga $a=b=0$.
Para $s<1$, integrar la serie geométrica en compactos y usar
$\int_0^s t^{2k}\log(s/t)\,dt=s^{2k+1}/(2k+1)^2$
recupera la serie de \eqref{md:compos:momento-catalan}.
Su convergencia absoluta dominada permite $s\uparrow1$ y restituye Catalán.
\end{proof}
La condición de frontera fija las dos constantes de integración.
La restitución utiliza la función constitutiva con su dominio; el
valor extremo de Catalán y los dos valores de canal son lecturas de
esa función ya construida. En el extremo, la identidad doblemente
diferenciada conserva su interpretación por límite de Abel.

El operador constitutivo $V=r_+P_++r_-P_-$ determina
\begin{equation}
 \widehat\varepsilon=r_+^2,\qquad \widehat\mu=r_-^2,
 \qquad \widehat Z=\frac{r_-}{r_+},\qquad
 \widehat c=(r_+r_-)^{-1}.
 \label{md:compos:lectores-constitutivos}
\end{equation}
Con $\psi=f^{-1}$ y
\begin{equation}
 \mathcal H_B(s)=\frac{12s^3}{1-s^9}
 -\frac{s^4}{1-s^{12}}-\frac{(\log s)^2}{20\pi},
 \label{md:compos:h-barbero}
\end{equation}
la factorización construida toma la forma
\begin{equation}
 \gamma=\mathcal H_B(\psi(r_-))-\mathcal H_B(\psi(r_+)).
 \label{md:compos:factorizacion-barbero}
\end{equation}
La parte logarítmica restituye exactamente la contribución angular de
Barbero, conservada en grados o transportada íntegramente a radianes.

\subsection{Restitución mediante velocidad normalizada y Barbero}
\label{md:compos:difeomorfismo}

\paragraph{Proposición.}
En la colección constitutiva normalizada de dos canales etiquetados, el mapa
\begin{equation}
 (r_+,r_-)\longmapsto(\widehat c,\gamma),\qquad
 (3/4,1)^2\longrightarrow(1,16/9)\times\mathbb R
 \label{md:compos:mapa-global}
\end{equation}
es un difeomorfismo. Su inversa recupera los dos autovalores etiquetados.
Los proyectores de hoja pertenecen a la representación de partida: dos
escalares no recuperan una base espacial arbitraria.

\paragraph{Demostración.}
La expresión regular
$f(s)=(1+s+s^2)/(1+s+s^2+s^3)$ da
\begin{equation}
 f'(s)=-\frac{s^2(s^2+2s+3)}{(1+s+s^2+s^3)^2}<0.
 \label{md:compos:derivada-f}
\end{equation}
Así, $f$ es analítica y estrictamente decreciente, con $f(0^+)=1$ y
$f(1^-)=3/4$. Para $0<s<1$,
\begin{equation}
 \mathcal H_B'(s)=
 \frac{36s^2(1+2s^9)}{(1-s^9)^2}
 -\frac{4s^3(1+2s^{12})}{(1-s^{12})^2}
 -\frac{\log s}{10\pi s}>0.
 \label{md:compos:derivada-h}
\end{equation}
En efecto, el cociente del primer término positivo entre el segundo término
positivo antes de restarlo es
\begin{equation}
 \frac{9}{s}\frac{1+2s^9}{1+2s^{12}}
 \left(\frac{1-s^{12}}{1-s^9}\right)^2>9.
 \label{md:compos:cota-derivada-h}
\end{equation}
El último sumando de \eqref{md:compos:derivada-h} también es positivo.
Por tanto, $F_B=\mathcal H_B\circ\psi$ tiene derivada estrictamente negativa,
$F_B(3/4^+)=+\infty$ y $F_B(1^-)=-\infty$.

Fijado $c=\widehat c$, escribimos $z=\log\widehat Z$ y
$r_\pm=c^{-1/2}\exp(\mp z/2)$. El intervalo admisible es
\begin{equation}
 |z|<a(c):=\min\left\{\log c,\log\frac{16}{9c}\right\}.
 \label{md:compos:dominio-z}
\end{equation}
La función $\Gamma_c(z)=F_B(r_-)-F_B(r_+)$ satisface
\begin{equation}
 \Gamma_c'(z)=\frac12
 \bigl[r_-F_B'(r_-)+r_+F_B'(r_+)\bigr]<0.
 \label{md:compos:derivada-gamma}
\end{equation}
Cuando $z$ aumenta hasta $a(c)$, o $r_-$ alcanza $1$, o $r_+$ alcanza $3/4$,
o suceden ambas cosas; $\Gamma_c$ tiende a $-\infty$. En el extremo opuesto
tiende a $+\infty$. Existe, por tanto, una única solución para cada $\gamma$
real. La derivada no nula y el teorema de la función implícita proporcionan
una inversa suave global. Esto prueba el enunciado. En la orientación
original $y>0$, se tiene $z<0$ y $\gamma>0$.

La velocidad normalizada conserva el producto de las respuestas; Barbero
distingue de forma monótona la distribución orientada entre ellas. Juntas
restituyen la información espectral que cada lectura aislada no determina.

El enunciado concierne a la colección de lectores constitutivos. La imagen
efectivamente producida por el generador HMT completo es un subconjunto de
esa colección: restringir a dicha imagen conserva la inyectividad y el orden
causal. La normalización y las marcas de hoja son datos necesarios de la
restitución. La prueba no introduce valores CODATA para elegir canales ni
afirma que cada punto sintético de la colección sea un estado generado.

\subsection{Restitución sobre la imagen generada}
\label{md:compos:restitucion-completa}

Una vez recuperados $r_\pm$ desde $(\widehat c,\gamma)$, se obtiene
$s_\pm=\psi(r_\pm)$. Entonces
\begin{equation}
 A_{\deg}=-\frac{3}{\pi}\log(s_+s_-),\qquad
 C^*_{\deg}=\frac{3}{\pi}\log\frac{s_-}{s_+}.
 \label{md:compos:restitucion-angular}
\end{equation}
Sobre la imagen generada y en su orientación positiva,
\begin{equation}
 \alpha=\frac{A_{\deg}}{1000},\qquad
 \eta_{\mathrm{ret}}=\frac{C^*_{\deg}}{2}-\alpha>0.
 \label{md:compos:restitucion-alpha-eta}
\end{equation}
Con las coordenadas regionales $\pi,\varphi,e$ y las secciones dimensionales
ya construidas, esta restitución devuelve $H_5$, el factor
$R_{\mathrm{act}}=H_5/\eta_{\mathrm{ret}}$ y la lectura electrónica completa
de la sección~\ref{md:compos:electron}. También devuelve la sección de
Planck. Si se conservan además la base posicional de doce bloques, el origen
y el acarreo, la identidad
\begin{equation}
 \kappa_{\mathrm{ag}}=\pi+e-\varphi-4-\alpha
 \label{md:compos:restitucion-k}
\end{equation}
permite la restitución del registro $K$ descrita por su lector. Se trata de
un recorrido de reconstrucción; su sentido inverso no constituye una nueva
flecha generadora desde valores físicos externos.

La gravedad y la conversión térmica conservan las secciones del mismo
estado. El teorema~\ref{md:rigidez:teorema} y
\eqref{md:rigidez:reloj} determinan primero $L_\alpha$ y su
representación circular $R_{108}=L_\alpha$:
\begin{equation}
 G_a=\frac{c_{\mathrm{int}}^3R_{108}^2}{\hbar_a},\qquad
 \frac{G_a\hbar_a}{c_{\mathrm{int}}^3}=R_{108}^2,
 \qquad
 k_B\Theta_{\mathrm{clk},a}\log3=\frac{h_a}{108t_0}.
 \label{md:compos:gravedad-termica}
\end{equation}
El radio cronológico, la velocidad física y la base térmica conservan sus
constructores. En particular, el reloj de la representación circular queda
ligado a $L_\alpha$ y $c_{\mathrm{int}}$. La primera
composición explicita la reciprocidad acción--gravedad; la segunda reúne
acción, duración e información trítica, con individualización de temperatura
y $k_B$ en sus bases declaradas.

\subsection{Definiciones estructurales de las lecturas compuestas}
\label{md:compos:definiciones}

\begin{description}
 \item[$\alpha$.] Coordenada escalar de clausura dodecafásica compatible con
 las regiones y con el registro aritmético--geométrico $K$; transmite esa
 compatibilidad al sector angular y de acción.
 \item[Contraángulo.] Coordenada angular de la clausura bidireccional que
 combina el coeficiente de la sección de acción de retorno con $\alpha$,
 en la coordenatización de $360$ unidades.
 \item[Planck reducida.] Sección de acción cuyo coeficiente regional y escala
 determinantal se construyen antes del contraángulo; en la elipse, medida de
 área orientada por unidad de fase paramétrica.
 \item[Electrón.] Sección central persistente; su lectura escalar compone
 incompatibilidad aditiva--multiplicativa, transferencia regional, defecto
 de clausura y retorno de acción. La ecuación completa conserva la sección,
 además de su valor de masa.
 \item[Barbero--Immirzi.] Funcional orientado del transporte angular y de los
 retornos de la incidencia hexada--octada; en la colección constitutiva,
 coordenada que junto con la velocidad normalizada restituye el espectro
 etiquetado.
 \item[Catalán.] Valor extremo del momento orientado de profundidad dos del
 cuarto de giro dodecafásico. Su colección se enlaza diferencialmente con la
 respuesta del vacío.
 \item[$G$.] Coeficiente común de la realización radial que relaciona
 energía y longitud conjugada conservando simultáneamente
 $H\Lambda=\hbar cI$ y $R\Lambda=L_\alpha^2I$.
 La composición longitudinal y polar determina su normalización;
 su expresión dimensional es $c^3L_\alpha^2/\hbar$.
 \item[Boltzmann.] Transductor entre la recta térmica y la energía de
 decisión, con normalización informacional trítica y base térmica declaradas.
\end{description}
Estas formulaciones intrínsecas se basan en los objetos y ecuaciones
construidos. Su alcance es el de esta realización; las identificaciones
físicas conservan sus demostraciones correspondientes y no se deducen de
la sola denominación estructural.

\subsection{Transporte del bloque electrónico nilpotente}
\label{md:compos:nilpotentes}

\paragraph{Proposición.}
Sean
\begin{equation}
 S=\begin{pmatrix}1&0\\0&-1\end{pmatrix},\qquad
 J=\begin{pmatrix}0&1\\-1&0\end{pmatrix},\qquad
 n_\pm=\frac{S\pm J}{2},\qquad
 D=\operatorname{diag}(r,r^{-1}),\quad r>0.
 \label{md:compos:bloque-nilpotente}
\end{equation}
Defínanse $n_{\pm,D}=Dn_\pm D^{-1}$ y $g=D^{-T}D^{-1}$. Entonces
\begin{equation}
 n_{\pm,D}^2=0,\qquad
 \{n_{+,D},n_{-,D}\}=I,\qquad
 (n_{+,D})^{\dagger_g}=n_{-,D},
 \label{md:compos:identidades-nilpotentes}
\end{equation}
donde $A^{\dagger_g}=g^{-1}A^Tg$. La elipse transportada desde
$Y^TY=2\hbar$ mediante $X=DY$ satisface $X^TgX=2\hbar$.

\paragraph{Demostración.}
El cálculo matricial directo da
$S^2=I$, $J^2=-I$ y $SJ+JS=0$. Por ello,
\[
 n_\pm^2=\tfrac14(S^2+J^2\pm SJ\pm JS)=0,
 \qquad
 n_+n_-+n_-n_+=\tfrac12(S^2-J^2)=I.
\]
Además, $S^T=S$ y $J^T=-J$ implican $n_+^T=n_-$. La conjugación por $D$
preserva productos, cuadrados y la identidad. Para el adjunto métrico,
\[
 g^{-1}(Dn_+D^{-1})^Tg
 =DD^T D^{-T}n_-D^T D^{-T}D^{-1}
 =Dn_-D^{-1}=n_{-,D}.
\]
Finalmente, $D^TgD=I$, de donde
$X^TgX=Y^TD^TgDY=Y^TY=2\hbar$. Quedan probadas las afirmaciones.

La comprobación algebraica exacta tiene lugar en
$\mathbb Q[r,r^{-1}]$; la restricción $r>0$ fija la realización elíptica
orientada utilizada. Esta composición conserva las funciones particulares
de los distintos factores $1/2$ del desarrollo: una coincidencia numérica
entre ellos no constituye una identificación de sus operadores.
